\documentclass[aspectratio=169, 11pt]{beamer}

% --- Theme ---
\usetheme{Boadilla}
\usecolortheme{beaver}

% --- Packages ---
\usepackage[T1]{fontenc}
\usepackage{textcomp}
\usepackage{graphicx}
\usepackage{amsmath, amssymb}
\newcommand{\Ang}{\text{\AA}}
\usepackage{tikz}
\usetikzlibrary{arrows.meta, positioning, shapes.geometric}
\usepackage{booktabs}

% --- Extra colors for TikZ diagrams ---
\definecolor{darkgray}{RGB}{80,80,80}
\definecolor{lightgray2}{RGB}{230,230,230}
\definecolor{deepred}{RGB}{140,20,20}

% --- Gray accents ---
\setbeamercolor{section in head/foot}{bg=gray!30, fg=darkred}
\setbeamercolor{subsection in head/foot}{bg=gray!15, fg=darkred}
\setbeamercolor{block body}{bg=gray!10}

% --- Top-right logo ---
\addtobeamertemplate{frametitle}{}{%
  \begin{tikzpicture}[remember picture, overlay]
    \node[anchor=north east, yshift=-2pt, xshift=-4pt]
      at (current page.north east) {\includegraphics[height=0.7cm]{metu.png}};
  \end{tikzpicture}%
}

% --- Title ---
\title[MEAM Mock Pitch]{Automated Construction of Multi-Element\\MEAM Interatomic Potentials}
\subtitle{DFT-Informed Neural Network Optimization}
\author{Can Durmaz}
\institute{METU}
\date{PHYS\,400 --- Mock Presentation}
\titlegraphic{\includegraphics[height=1.5cm]{metu.png}}

\begin{document}

% --- Slide 1: Title ---
\begin{frame}[plain]
\titlepage
\end{frame}

% --- Slide 2: The Problem ---
\begin{frame}{The Problem}
\begin{itemize}
  \item Industrial alloys (Al\,7075, steels) are chosen by their \textbf{elastic properties} --- yet measuring them experimentally is \textbf{slow and expensive}.
  \item Molecular dynamics can predict these properties \emph{computationally}, but only with accurate \textbf{interatomic potentials}.
  \item The Modified Embedded Atom Method (MEAM) is the standard for metallic alloys.
\end{itemize}

\vspace{1em}
\begin{block}{Goal}
Build an automated pipeline that generates MEAM potentials for \emph{arbitrary} multi-element alloys.
\end{block}
\end{frame}

% --- Slide 3: Why Existing Approaches Fall Short ---
\begin{frame}{Why Existing Approaches Fall Short}
Published MEAM files cover only \textbf{subsets} of elements. No single file spans Al\,7075 (Al, Zn, Mg, Cu, Cr, Mn, Fe, Si, Ti).

\vspace{0.5em}
\centering
\small
\begin{tabular}{l|cccccccccccc}
\toprule
\textbf{Library File} & Al & Co & Cr & Cu & Fe & Mg & Mn & Mo & Ni & Si & Ti & Zn \\
\midrule
AlMgZn           & \checkmark &   &   &   &   & \checkmark &   &   &   &   &   & \checkmark \\
AlSiCuMgFe       & \checkmark &   &   & \checkmark & \checkmark & \checkmark &   &   &   & \checkmark &   &   \\
CuMo             &   &   &   & \checkmark &   &   &   & \checkmark &   &   &   &   \\
FeMnNiTiCuCrCoAl & \checkmark & \checkmark & \checkmark & \checkmark & \checkmark &   & \checkmark &   & \checkmark &   & \checkmark &   \\
\bottomrule
\end{tabular}

\vspace{0.7em}
\raggedright
\begin{itemize}\small
  \item \textbf{Cross-interaction parameters} between files are undefined $\Rightarrow$ cannot simulate combined compositions.
  \item \textbf{Pure ML potentials} (MLIPs) are accurate but demand massive datasets and sacrifice MEAM's transparency.
\end{itemize}
\end{frame}

% --- Slide 4: Pipeline Overview (interim-style, 5 stages + ML branch) ---
\begin{frame}{Pipeline Overview}
\centering
\begin{tikzpicture}[
    node distance=0.5cm and 0.35cm,
    stage/.style={rectangle, rounded corners, draw=darkred, fill=darkred!10,
                  text width=2.5cm, minimum height=1.0cm, align=center,
                  font=\scriptsize\bfseries},
    planned/.style={rectangle, rounded corners, draw=darkred, dashed, fill=darkred!3,
                  text width=2.5cm, minimum height=1.0cm, align=center,
                  font=\scriptsize\bfseries\itshape},
    aux/.style={rectangle, rounded corners, draw=blue!55!black, dashed, fill=blue!6,
                  text width=2.5cm, minimum height=0.9cm, align=center,
                  font=\scriptsize\bfseries},
    arrow/.style={-{Stealth[length=2.5mm]}, thick, darkred},
    plannedarrow/.style={-{Stealth[length=2.5mm]}, thick, dashed, darkred},
    auxarrow/.style={-{Stealth[length=2.5mm]}, thick, dashed, blue!55!black},
    label/.style={font=\tiny, text=darkred!70!black, align=center},
    auxlabel/.style={font=\tiny, text=blue!55!black, align=center},
  ]
  \node[stage] (s1) {Stage 1\\Configurations};
  \node[stage, right=of s1] (s2) {Stage 2\\MD: $E,\nu,C_{ij}$};
  \node[stage, right=of s2] (s3) {Stage 3\\DFT};
  \node[stage, right=of s3] (s4) {Stage 4\\MEAM init};
  \node[planned, right=of s4] (s5) {Stage 5\\NN optim\\(planned)};

  \draw[arrow] (s1) -- (s2);
  \draw[arrow] (s2) -- (s3);
  \draw[arrow] (s3) -- (s4);
  \draw[plannedarrow] (s4) -- (s5);

  \node[label, below=0.15cm of s1] {composition\\records};
  \node[label, below=0.15cm of s2] {elastic\\database};
  \node[label, below=0.15cm of s3] {DFT\\reference};
  \node[label, below=0.15cm of s4] {merged MEAM\\library};
  \node[label, below=0.15cm of s5] {refitted MEAM\\parameters};

  % Divergent ML sanity-check branch (does not feed back into main flow)
  \node[aux, above=1.0cm of s3] (ml) {ML branch\\(sanity check)};
  \draw[auxarrow] (s2.north) |- (ml.west);
  \node[auxlabel, right=0.15cm of ml] {$(E,\nu)$\\predictor};
\end{tikzpicture}
\end{frame}

% --- Slide 5: Pipeline Stages at a Glance (interim-style: # / Description / Framework / Status) ---
\begin{frame}{Pipeline Stages at a Glance}
\centering
\small
\begin{tabular}{clll}
\toprule
\textbf{Stage} & \textbf{Description} & \textbf{Framework} & \textbf{Status} \\
\midrule
1 & Random alloy configuration generation   & Python                  & Complete \\
2 & MD elastic-property evaluation          & LAMMPS Python API       & Complete \\
3 & DFT reference calculations              & Quantum Espresso + ASE  & Complete \\
4 & MEAM initialization \& library merging  & Python                  & Complete \\
5 & Neural-network surrogate optimization   & TensorFlow + LAMMPS     & \emph{Planned} \\
\bottomrule
\end{tabular}

\vspace{1em}
\begin{block}{Key Idea}
Use neural networks as a \textbf{surrogate} that learns the MEAM-parameters $\to$ elastic-properties map, then inverse-optimize parameters to hit target $(E, \nu)$.
\end{block}
\end{frame}

\end{document}
